\documentclass{article}
\usepackage[UTF8]{ctex}
\usepackage{amsmath,amssymb,amsfonts}
\usepackage{enumitem}
\usepackage{geometry}
\geometry{a4paper, margin=1in}
% === 错题本排版优化（不参与题目标识） ===
\setlength{\parindent}{0pt}
\setlength{\parskip}{0.35em plus 0.10em minus 0.08em}
\setlength{\emergencystretch}{3em}
\setlist{itemsep=0.25em, topsep=0.40em, parsep=0pt}
\allowdisplaybreaks[2]
% === 错题本排版优化结束 ===
\begin{document}
\section*{西南交通大学2025年数学分析真题}
\section*{171．西南交通大学2025年数学分析试卷}

\subsection{题目1}
1．（10 分）证明： \(\lim _{n \rightarrow \infty} \sin n\) 不存在．

\subsection{题目2}
2．（10 分）设 \(f(x)\) 是 \(\mathbb{R}\) 上周期为 \(T\) 的连续函数．证明： \(\lim _{t \rightarrow+\infty} \frac{1}{t} \int_{0}^{t} f(x) \mathrm{d} x=\frac{1}{T} \int_{0}^{T} f(x) \mathrm{d} x\) ．

\subsection{题目3}
3．（10 分）设函数 \(f(x)\) 在 \([1,+\infty)\) 上连续可微，且 \(\int_{1}^{+\infty} f(x) \mathrm{d} x\) 与 \(\int_{1}^{+\infty} f^{\prime}(x) \mathrm{d} x\) 均收敛，证明：

\[
\lim _{x \rightarrow+\infty} f(x)=0
\]

\subsection{题目4}
4．（15 分）设函数 \(f(x)\) 在 \((a, b)\) 内处处有导数 \(f^{\prime}(x)\) 。证明：\((a, b)\) 中的点或者为 \(f^{\prime}(x)\) 的连续点，或者为 \(f^{\prime}(x)\) 的第二类间断点．

\subsection{题目5}
5．（15 分）设函数 \(f(x)\) 在 \([1,+\infty)\) 上连续，在 \((1,+\infty)\) 上可导，已知函数 \(e^{-x} f^{\prime}(x)\) 在 \((1,+\infty)\) 上有界，证明：函数 \(e^{-x} f(x)\) 在 \((1,+\infty)\) 上也有界。

\subsection{题目6}
6．（10 分）求极限 \(\lim _{x \rightarrow 0} \frac{\ln \left(1+\sin ^{2} x\right)-6(\sqrt[3]{2-\cos x}-1)}{x^{4}}\) ．

\subsection{题目7}
7．（10 分）证明：有限闭区间上的单调函数必定可积．

\subsection{题目8}
8．（10 分）求幂级数 \(\sum_{n=0}^{\infty} \frac{n^{2}+1}{2^{n} n!} x^{n}\) 的和函数．

\subsection{题目9}
9．（15 分）\(K\) 为 \(\mathbb{R}^{n}\) 中的紧集，\(f: K \rightarrow \mathbb{R}^{n}\) 为连续映射，证明：\(f\) 在 \(K\) 上一致连续．

\subsection{题目10}
10．（15 分）设 \(u(x, y)\) 所有二阶偏导数都连续，且

\[
\frac{\partial^{2} u}{\partial x^{2}}-\frac{\partial^{2} u}{\partial y^{2}}=0, u(x, 2 x)=x, u_{x}(x, 2 x)=x^{2}
\]

求 \(u_{x x}(x, 2 x), u_{x y}(x, 2 x), u_{y y}(x, 2 x)\) ．

\begin{itemize}

\subsection{题目11}
\end{itemize}

\subsection{题目11}
原题边界“11．”：（15 分）计算积分 \(I=\iiint_{V} \frac{\mathrm{~d} V}{\rho^{2}}, \rho\) 为点 \((x, y, z)\) 到 \(x\) 轴的距离，\(V\) 表示一个棱台，它的顶点坐标为
\[
A(0,0,1), B(0,1,1), C(1,1,1), D(0,0,2), E(0,2,2), F(2,2,2) .
\]

\subsection{题目12}

\subsection{题目12}
原题边界“12．”：（15 分）求 \(\iint_{\sum} x \mathrm{~d} y \mathrm{~d} z+y \mathrm{~d} z \mathrm{~d} x+z \mathrm{~d} x \mathrm{~d} y, \sum\) 为上半球面 \(z=\sqrt{R^{2}-x^{2}-y^{2}}, R>0\) ，方向取上侧．
\end{document}
